\Needspace{9\baselineskip}\begin{center}\itshape \textit{Application of the preceding investigations to trigonometric functions.}\\\textit{A method of identifying the angles corresponding to the individual roots \(\Omega\).}\end{center}


\Needspace{6\baselineskip}\Article{361}
It remains for us to examine more closely the connection between the roots \(\Omega\) and the trigonometric functions of the angles \(\frac Pn,\allowbreak \frac{2P}n,\allowbreak \frac{3P}n\ldots\frac{(n-1)P}n\). The method we have given for finding the roots \(\Omega\) still leaves uncertain (unless sine tables have been consulted during the work in the way described above, which would, however, be less direct) \textit{which} roots correspond to \textit{each} of these angles, \(\text{that is}\) which root is \(=\cos\frac Pn+i\sin\frac Pn\), which is \(=\cos\frac{2P}n+i\sin\frac{2P}n\), etc. This uncertainty is easily removed by considering that the cosines of the angles \(\frac Pn,\allowbreak \frac{2P}n,\allowbreak \frac{3P}n\ldots\frac{(n-1)P}{2n}\) steadily decrease (provided their signs are also taken into account), while their sines are all positive; the angles \(\frac{(n-1)P}n,\allowbreak \frac{(n-2)P}n,\allowbreak \frac{(n-3)P}n\ldots\frac{(n+1)P}{2n}\), on the other hand, have respectively the same cosines as those angles, but negative sines whose absolute magnitudes equal the sines of those angles. Thus the two roots among \(\Omega\) with the greatest real parts (equal to one another) will correspond to the angles \(\frac Pn,\allowbreak \frac{(n-1)P}n\): to the former, the root in which the imaginary quantity \(i\) is multiplied by a positive quantity; to the latter, the root in which \(i\) is multiplied by a negative quantity. Of the \(n-3\) remaining roots, those with the greatest real parts will again correspond to the angles \(\frac{2P}n,\allowbreak \frac{(n-2)P}n\), and so on. — Once the root corresponding to the angle \(\frac Pn\) has been identified, those corresponding to the remaining angles can also be distinguished by the fact that, if the former is supposed to be \(=[\lambda]\), the angles \(\frac{2P}n,\allowbreak \frac{3P}n,\allowbreak \frac{4P}n\), etc., evidently correspond to the roots \([2\lambda],\allowbreak [3\lambda],\allowbreak [4\lambda]\), etc. Thus in the example of Article 353 one sees at once that the angle \(\frac1{19}P\) can correspond to no root other than \([11]\), and the angle \(\frac{18}{19}P\) to the root \([8]\); similarly, the angles \(\frac2{19}P,\allowbreak \frac{17}{19}P,\allowbreak \frac3{19}P,\allowbreak \frac{16}{19}P\), etc., correspond to the roots \([3],\allowbreak [16],\allowbreak [14],\allowbreak [5]\), etc. In the example of Article 354, the angle \(\frac1{17}P\) evidently corresponds to the root \([1]\), the angle \(\frac2{17}P\) to \([2]\), etc. In this way, therefore, the cosines and sines of the angles \(\frac Pn,\allowbreak \frac{2P}n\), etc., will be fully determined.
\SourceLeaf{460}

\clearpage
\begingroup\abovedisplayskip=6pt plus2pt minus1pt \belowdisplayskip=6pt plus2pt minus1pt \abovedisplayshortskip=0pt plus1pt \belowdisplayshortskip=3pt plus1pt minus1pt
\Needspace{9\baselineskip}\begin{center}\itshape \textit{Tangents, cotangents, secants and cosecants are derived from sines and cosines without division.}\end{center}


\Needspace{6\baselineskip}\Article{362}
As for the remaining trigonometric functions of these angles, they could easily be derived from the corresponding cosines and sines by the usual methods: namely, the secants and tangents by dividing unity and the sines by the cosines, and the cosecants and cotangents by dividing unity and the cosines by the sines. But the same result is usually obtained more conveniently with the help of the following formulas, without divisions, by additions alone. Let \(\omega\) be any of the angles \(\frac Pn,\allowbreak \frac{2P}n\ldots\frac{(n-1)P}n\), and let \(\cos\omega+i\sin\omega=R\), so that \(R\) will be one of the roots \(\Omega\), and \[\begin{aligned}\cos\omega&=\tfrac12(R+\tfrac1R)=\frac{1+RR}{2R},\\\sin\omega&=\tfrac1{2i}(R-\tfrac1R)=\frac{i(1-RR)}{2R}\end{aligned}\] Hence \[\begin{gathered}\sec\omega=\frac{2R}{1+RR}.\quad\operatorname{tang}\omega=\frac{i(1-RR)}{1+RR},\\\operatorname{cosec}\omega=\frac{2Ri}{RR-1},\quad\operatorname{cotang}\omega=\frac{i(RR+1)}{RR-1}\end{gathered}\] We will now show how to transform the numerators of these four fractions so that they become divisible by their denominators.

I. Since \(R=R^{n+1}=R^{2n+1}\), we have \(2R=R+R^{2n+1}\), an expression evidently divisible by \(1+RR\), since \(n\) is odd. Hence \[\sec\omega=R-R^3+R^5-R^7\ldots+R^{2n-1}\] and therefore (since \(\sin\omega=-\sin(2n-1)\omega,\allowbreak \ \sin3\omega=-\sin(2n-3)\omega\), etc., evidently imply \(\sin\omega-\sin3\omega+\sin5\omega\ldots+\sin(2n-1)\omega=0\)) \[\sec\omega=\cos\omega-\cos3\omega+\cos5\omega\ldots+\cos(2n-1)\omega\] or finally (since \(\cos\omega=\cos(2n-1)\omega,\allowbreak \ \cos3\omega=\cos(2n-3)\omega\), etc.), \[=2(\cos\omega-\cos3\omega+\cos5\omega\ldots\mp\cos(n-2)\omega)\pm\cos n\omega\] the upper or lower sign applying according as \(n\) has the form \(4k+1\) or \(4k+3\). This formula can evidently also be written as \[\sec\omega=\pm(1-2\cos2\omega+2\cos4\omega\ldots\pm2\cos(n-1)\omega)\]

II. Similarly, substituting \(1-R^{2n+2}\) for \(1-RR\) gives
\SourceLeaf{461}
\[\operatorname{tang}\omega=i(1-RR+R^4-R^6\ldots-R^{2n})\] or (since \(1-R^{2n}=0,\allowbreak \ RR-R^{2n-2}=2i\sin2\omega,\allowbreak \ R^4-R^{2n-4}=2i\sin4\omega\), etc.), \[\operatorname{tang}\omega=2(\sin2\omega-\sin4\omega+\sin6\omega\ldots\mp\sin(n-1)\omega)\]

III. Since we have \(1+RR+R^4\ldots+R^{2n-2}=0\), it follows that \[\begin{aligned}n&=n-1-RR-R^4\ldots R^{2n-2}\\&=(1-1)+(1-RR)+(1-R^4)\ldots+(1-R^{2n-2})\end{aligned}\] and each part of this sum is divisible by \(1-RR\). Hence \[\begin{aligned}\frac n{1-RR}&=1+(1+RR)+(1+RR+R^4)\ldots\\&\quad+(1+RR+R^4\ldots+R^{2n-4})\\&=(n-1)+(n-2)RR+(n-3)R^4\ldots+R^{2n-4}\end{aligned}\] so that, multiplying by \(2\), subtracting \[0=(n-1)(1+RR+R^4\ldots+R^{2n-2})\] and again multiplying by \(R\), we obtain \[\begin{aligned}\frac{2nR}{1-RR}&=(n-1)R+(n-3)R^3+(n-5)R^5\ldots\\&\quad-(n-3)R^{2n-3}-(n-1)R^{2n-1}\end{aligned}\] from which we immediately deduce \[\begin{aligned}\operatorname{cosec}\omega&=\frac1n\bigl((n-1)\sin\omega+(n-3)\sin3\omega\ldots\\&\qquad\qquad-(n-1)\sin(2n-1)\omega\bigr)\\&=\frac2n\bigl((n-1)\sin\omega+(n-3)\sin3\omega\\&\qquad\qquad+\text{etc.}+2\sin(n-2)\omega\bigr)\end{aligned}\] This formula can also be written as \[\operatorname{cosec}\omega=-\frac2n(2\sin2\omega+4\sin4\omega+6\sin6\omega\ldots+(n-1)\sin(n-1)\omega)\]

IV. Multiplying the value of \(\frac n{1-RR}\) given above by \(1+RR\) and subtracting \[0=(n-1)(1+RR+R^4\ldots+R^{2n-2})\] gives \[\frac{n(1+RR)}{1-RR}=(n-2)RR+(n-4)R^4+(n-6)R^6\ldots-(n-2)R^{2n-2}\]
\SourceLeaf{462}
from which it immediately follows that \[\begin{aligned}\operatorname{cotang}\omega&=\frac1n\bigl((n-2)\sin2\omega+(n-4)\sin4\omega\\&\qquad+(n-6)\sin6\omega\ldots-(n-2)\sin(n-2)\omega\\&=\frac2n((n-2)\sin2\omega+(n-4)\sin4\omega\ldots\\&\qquad+3\sin(n-3)\omega+\sin(n-1)\omega)\end{aligned}\] This formula may also be written in the following way: \[\operatorname{cotang}\omega=-\frac2n(\sin\omega+3\sin3\omega\ldots+(n-2)\sin(n-2)\omega)\]\endgroup


